
\sgaStableTarget{sga6.e10.sec.6.3}\subsection*{6.3.}
Let us first show how 6.2.3 follows from 6.2.2. For this, it plainly suffices to prove (6.2.5) for the
@@M0@@ @@M1@@, so one may suppose @@M2@@, @@M3@@. By 6.2.2, @@M4@@ factors through @@M5@@, i.e., there is a commutative diagram
@@M6@@
where @@M7@@ is the canonical epimorphism and @@M8@@ is a representation of @@M9@@ in @@M10@@. But the homomorphism
@@M11@@
has a section, so that the monomorphism
@@M12@@
induced by det is an isomorphism. On the other hand, it is known (SGA 3 VIII 1.3) that every representation of @@M13@@ in itself is of the form
@@M14@@
where @@M15@@ is a section of @@M16@@, whence the assertion.

Corollary 6.2.6 is a simple reformulation of 6.2.3. Note that 6.2.3 also follows immediately from SGA 3 XXII 6.2.1, but we wish to give an elementary proof here via 6.2.2.
